% ---
% Reemplaza líneas 2941–2996 de Docs/Technical_Report.tex.
% Decisiones que lo respaldan: D-005, D-012, H-001.
% Capítulo 4 desde el 2026-10-14 (antes §5).
% ---

\subsection{Validación y limpieza de los datos}

% fuente: H-001, D-005
La carga y unificación de los cuatro archivos CSV produjo
un DataFrame consolidado de 3\,568 registros, con las
columnas \texttt{finding\_type} (masa / calcificación) y
\texttt{split} (train / test) agregadas para facilitar
el análisis por subconjunto. La Tabla~\ref{tab:eda_resumen}
presenta el resumen de la composición del dataset.

% fuente: H-001, D-005, H-038
\begin{table}[htbp]
\centering
\caption{Composición del dataset CBIS-DDSM unificado}
\label{tab:eda_resumen}
\renewcommand{\arraystretch}{1.3}
\begin{tabularx}{\textwidth}{@{}Xrr@{}}
\toprule
\textbf{Subconjunto} &
\textbf{Train} & \textbf{Test} \\
\midrule
Masas (mass)           & 1\,318 & 378 \\
Calcificaciones (calc) & 1\,546 & 326 \\
\midrule
\textbf{Total}         & \textbf{2\,864} &
                         \textbf{704} \\
\bottomrule
\end{tabularx}
\end{table}

% fuente: D-005, H-001; 1_EDA.ipynb
La etiqueta objetivo \texttt{pathology} presentó tres
valores distintos en los datos originales:
\texttt{BENIGN} (1\,548 casos),
\texttt{MALIGNANT} (1\,457 casos) y
\texttt{BENIGN\_WITHOUT\_CALLBACK} (563 casos).
Los casos \texttt{BENIGN\_WITHOUT\_CALLBACK} fueron
normalizados a \texttt{BENIGN} por equivalencia clínica,
resultando en la codificación binaria
0\,=\,benigno (2\,111 casos, 59.2\,\%) y
1\,=\,maligno (1\,457 casos, 40.8\,\%).

La validación de nulos confirmó que las columnas
\texttt{mass\_shape}, \texttt{mass\_margins},
\texttt{calc\_type} y \texttt{calc\_distribution}
presentan valores faltantes de naturaleza estructural:
los registros de calcificaciones no tienen valores en
las columnas de masas y viceversa, por diseño del
dataset. No se identificaron errores de datos que
requirieran corrección.

% fuente: D-012, H-001
La verificación del enlace entre registros CSV y
archivos DICOM en disco arrojó una cobertura del
100\,\% (3\,568/3\,568 registros con archivo
DICOM encontrado), alcanzada tras recuperar las 84 series
faltantes mediante la API REST de TCIA; la descarga inicial
con el NBIA Data Retriever dejaba una cobertura del 97.6\,\%.
